%==============================================================================
% KATA PENGANTAR
%==============================================================================
\chapter*{Kata Pengantar}
\addcontentsline{toc}{chapter}{Kata Pengantar}

Puji syukur penulis panjatkan kepada Tuhan Yang Maha Esa atas rahmat-Nya penulis dapat menyelesaikan laporan mengenai topik ini dengan baik. Laporan ini disusun sebagai bentuk pertanggungjawaban akademik atas pelaksanaan Praktik Kerja Lapangan (PKL) mengenai topik ini yang telah dilakukan oleh penulis di PT. Bangga Teknologi Produktif (Karla Bionics).

Praktik Kerja Lapangan (PKL) ini merupakan salah satu bagian dari proses belajar yang memberikan kesempatan kepada mahasiswa untuk mempraktikkan ilmu yang telah diperoleh di perkuliahan. Melalui kegiatan ini penulis memperoleh pengalaman pengembangan perangkat lunak, memahami alur kerja industri, etika profesionalisme, serta pentingnya kolaborasi antara individu dan tim dalam bekerja.

Dalam pelaksanaannya, penulis ingin mengucapkan rasa terima kasih dengan segala kerendahan hati kepada:
\begin{enumerate}
    \item Ayah, Ibu, dan keluarga penulis yang senantiasa memberikan doa, dukungan, dan semangat dalam setiap langkah perjalanan penulis.
    \item Bapak Andre Febrianto, S.Kom., M.Eng., dosen pembimbing di Institut Teknologi Sumatera, yang telah memberikan arahan, saran, serta motivasi selama penyusunan laporan ini.
    \item Pihak PT Bangga Teknologi Produktif (Karla Bionic), khususnya Bapak Wildan Trusaji, S.T., M.T., dan Arief Indra Muharam selaku mentor lapangan yang telah memberikan kesempatan, bimbingan, kritik dan saran selama pelaksanaan penelitian.
    \item Rekan-rekan seperjuangan yang selalu memberikan dorongan moral serta berbagi kritik dan saran.
\end{enumerate}

Penulis menyadari laporan yang ditulis masih memiliki kekurangan. Oleh karena itu, penulis mengharapkan kritik dan saran yang membangun dari para pembaca. Harapan penulis adalah laporan ini dapat memberikan manfaat bagi pihak perusahaan, institusi pendidikan, serta pembaca yang tertarik dengan topik ini.
